
\begin{obereflection}
\textbf{Latihan 13.1}
1. Mengapa algoritma knapsack sederhana hanya dapat dipakai untuk enkripsi,
   tetapi tidak untuk dekripsi?
2. Apa yang terjadi jika $\text{PBB}(n, m) \neq 1$ pada pembangkitan kunci?
3. Mengapa algoritma \textit{greedy} dapat dipercaya pada barisan
   \textit{superincreasing}, tetapi tidak pada barisan normal?
4. Diberikan barisan $W = \{3, 7, 14, 29, 60, 118\}$ dan $M = 96$. Tentukan
   apakah barisan itu \textit{superincreasing}, lalu pecahkan $M$ dengan
   algoritma \textit{greedy}.
5. Diketahui $W_{priv} = \{1, 3, 7, 15, 31\}$ dengan $m = 71$ dan $n = 23$.
   Hitung $W_{pub}$ dan $n^{-1} \bmod 71$, lalu enkripsikan blok
   \code{10110}. Periksa bahwa dekripsinya mengembalikan blok yang sama.
6. Bila $m$ dipilih lebih kecil daripada jumlah seluruh elemen $W_{priv}$,
   langkah mana pada Subbagian~\ref{subsec:bukti-knapsack} yang gugur? Jelaskan.
7. Mengapa plainteks harus dipecah menjadi blok yang panjangnya sama dengan
   banyaknya elemen kunci publik?
\end{obereflection}

\begin{obereflection}
\textbf{Latihan 13.2}
1. Tunjukkan bahwa kunci publik $\{62, 93, 81, 88, 102, 37\}$ bukan barisan
   \textit{superincreasing}. Mengapa sifat itu hilang setelah transformasi
   modulo?
2. Mengapa $C = \sum_i b_i w'_i$ tidak perlu diambil modulo $m$, sedangkan
   dekripsi memakai $C \bmod m$?
3. Shamir memulihkan kunci privat Merkle--Hellman dalam waktu polinomial.
   Apakah hal itu berarti persoalan knapsack umum juga dapat dipecahkan dalam
   waktu polinomial? Jelaskan.
4. Sumber menganjurkan 250 elemen berukuran 200 sampai 400 bit, tetapi juga
   modulus 100 sampai 200 bit. Tunjukkan bahwa kedua anjuran itu tidak dapat
   dipenuhi sekaligus. Berapa panjang modulus yang diperlukan?
5. Mengapa penemuan $(n, m)$ oleh penyerang sudah cukup untuk membuka seluruh
   pesan, tanpa perlu memecahkan \textit{knapsack} publik?
\end{obereflection}
